\cleardoublepage

\chapter{Conclusiones y trabajos futuros}
\label{conclusion-y-trabajos-futuros}

Este capítulo cierra la memoria. Primero recoge las conclusiones del trabajo en relación con
los objetivos planteados al inicio, valorando en qué medida se han alcanzado y qué aporta cada
fase del sistema. A continuación reflexiona sobre las dificultades encontradas durante el
desarrollo, tanto las de naturaleza física como las de carácter práctico. Por último, propone
un conjunto de líneas de continuación ordenadas por ambición, que prolongan de forma natural lo
realizado.

\section{Conclusiones}
\label{conclusiones}

El objetivo general de este Trabajo de Fin de Máster era diseñar, implementar y validar un
\textit{pipeline} de aprendizaje profundo de extremo a extremo capaz de detectar eventos
sísmicos, modelar patrones de actividad y generar mapas de riesgo a corto plazo, con un nivel
de rigor metodológico comparable al de las publicaciones del área. A la luz de los resultados
del capítulo anterior, puede afirmarse que ese objetivo se ha cumplido: el sistema integra
cinco modelos en un flujo coherente y reproducible, y cada fase aporta un resultado claro y
verificable. Más allá de construir el sistema, el trabajo ha conseguido algo igualmente
valioso, que es \emph{medir con honestidad} qué parte del problema es abordable con los datos
disponibles y qué parte no, distinción que constituye la principal aportación científica de la
memoria. A continuación se detallan las conclusiones por objetivos parciales.

\paragraph{Construcción de una base experimental reproducible.} Se ha construido un conjunto
de datos coherente a partir de dos fuentes abiertas, STEAD para las formas de onda y el
catálogo del USGS para el pronóstico, con regiones, ventanas y umbrales justificados. Un
análisis empírico ha mostrado por qué cada fuente es adecuada para su tarea: STEAD aporta
trazas etiquetadas de alta calidad para detección, mientras que el catálogo USGS ofrece la
cobertura completa y homogénea que exige el cálculo de características de pronóstico, como
confirmó la comparación directa entre ambos catálogos. Todas las decisiones de preprocesado
--magnitud de completitud, aislamiento de eventos independientes mediante
\textit{declustering}, particiones estrictamente temporales y normalización ajustada solo en
entrenamiento-- se han tomado con criterios documentados y reproducibles. Este cuidado
metodológico, lejos de ser un detalle accesorio, es lo que permite confiar en las conclusiones
posteriores: sin él, cualquier resultado quedaría en entredicho por posibles fugas de
información.

\paragraph{Detección y \textit{picking} de fases.} Las arquitecturas preentrenadas PhaseNet y
EQTransformer alcanzan un rendimiento alto sobre datos no vistos, con un \textit{F1-score} de
hasta 0.93 en la fase P --que asciende por encima de 0.96 si se relaja ligeramente la
tolerancia temporal-- y errores de localización de pocas centésimas de segundo, comparables a
los de un analista humano. Esto confirma que la detección automática de fases es hoy una tarea
resuelta de forma fiable mediante aprendizaje profundo, y proporciona una base firme sobre la
que apoyar el resto del sistema. La elección de dos arquitecturas de familias distintas --una
\textit{U-Net} convolucional y una red multitarea con atención-- ha permitido además comprobar
que el resultado no depende de un diseño concreto.

\paragraph{Modelado de la dimensión temporal y espacial.} El trabajo ha contrastado de forma
sistemática un enfoque puramente temporal frente a uno espacial. El modelo temporal global no
supera a las referencias, lo que indica que la evolución temporal agregada de la sismicidad no
contiene, por sí sola, información suficiente para el pronóstico. La incorporación de la
dimensión espacial, en cambio, sí produce una mejora sustancial: tanto la red de grafos como el
Transformer geoespacial alcanzan un AUC en torno a 0.70--0.73, y el Transformer supera a la
climatología de Poisson de forma estadísticamente significativa. La representación del espacio
se revela, por tanto, como el factor decisivo, y el mecanismo de atención del Transformer
resulta ligeramente más eficaz que la propagación por grafo de la GNN. Este resultado tiene
valor práctico inmediato: orienta el diseño de futuros sistemas hacia las arquitecturas que
modelan explícitamente la estructura espacial de la sismicidad.

\paragraph{Localización de la información predictiva.} La aportación metodológica más relevante
del trabajo es la separación explícita de las dos preguntas que encierra el problema --dónde y
cuándo-- mediante una métrica diseñada al efecto, la AUC temporal por celda. Su aplicación
permite afirmar con precisión que la capacidad del sistema reside en la dimensión espacial: el
modelo identifica muy bien las zonas crónicamente peligrosas, mientras que la anticipación
temporal a 30 días, una vez aislada del componente espacial, coincide con la tasa base. Este
resultado, confirmado por varios estudios de ablación que lo mantienen estable frente a cambios
de arquitectura, de ventana y de preprocesado, es plenamente coherente con el consenso
científico actual sobre la dificultad del pronóstico a corto plazo. Constituye un hallazgo
honesto y bien fundamentado, más valioso que una cifra global que mezclase ambas dimensiones y
diese una impresión inflada del rendimiento. Saber con certeza dónde reside --y dónde no-- la
información aprovechable es, en sí mismo, un avance.

\paragraph{Alcance y generalización del modelo.} Un experimento de transferencia a nueve regiones
de los cinco continentes ha permitido delimitar con precisión el alcance del sistema. Aplicado sin
reentrenar fuera de California+Alaska, el modelo no generaliza: en todas las regiones queda por
debajo de su climatología local y su AUC medio (0.504) es indistinguible del azar. Este resultado,
lejos de restar valor al trabajo, lo precisa: confirma que lo que el sistema aprende es el
\emph{mapa espacial de peligro} de las zonas con las que se entrenó, una estructura propia de cada
contexto tectónico y no un patrón universal. La lectura del experimento es además coherente con la
física: el modelo transfiere algo mejor a las subducciones muy activas de Asia --semejantes a la
de Alaska-- y peor a las regiones de geometría difusa o sísmicamente tranquila. Queda así claro
que el sistema es una herramienta de cartografía \emph{regional}, que debe ajustarse a cada zona,
y no un pronosticador global.

\paragraph{Cuantificación de la incertidumbre y mapas de riesgo.} La capa final basada en un
proceso gaussiano cierra el \textit{pipeline} produciendo mapas continuos de riesgo acompañados
de su incertidumbre. La calibración previa de las probabilidades, validada de forma
independiente sobre una partición reservada, reduce el \textit{Brier score} en dos tercios, de
modo que los mapas expresan frecuencias de ocurrencia realistas y no valores inflados por la
ponderación de clase. La validación cuantitativa confirma que los mapas concentran el riesgo en
las zonas donde efectivamente ocurren los terremotos, con un AUC espacial de 0.80 y un riesgo
medio más del doble en las celdas con eventos reales. El mapa de incertidumbre, además, crece de
forma coherente en las zonas con menos datos, lo que comunica al usuario dónde debe confiar más
y dónde menos en la predicción. El sistema entrega, así, un producto cartográfico accionable y
honesto sobre su propia confianza.

\medskip
En conjunto, el trabajo demuestra que un \textit{pipeline} de aprendizaje profundo construido
con rigor es capaz de cartografiar el peligro sísmico espacial con calidad y de acotar con
precisión los límites de la predicción temporal a corto plazo. Esta combinación --un producto
útil y una delimitación honesta de lo que es y no es posible-- es, precisamente, lo que
distingue a un trabajo metodológicamente sólido de una mera optimización de cifras. El sistema
no solo responde a la pregunta que se le plantea, sino que sabe explicar con datos por qué
responde mejor a unas preguntas que a otras, lo que lo convierte en una herramienta fiable
sobre la que seguir construyendo.

\section{Limitaciones y dificultades}
\label{limitaciones}

El desarrollo del trabajo ha encontrado varias dificultades que conviene explicitar, tanto por
honestidad científica como porque orientan con claridad las líneas futuras.

La principal limitación es de naturaleza física y no metodológica: la información necesaria para
anticipar el \emph{momento} de un terremoto a 30 días no parece estar presente en el catálogo de
eventos. La acumulación de tensión que precede a un sismo se produce en escalas de tiempo de
años y de forma mayoritariamente silenciosa, por lo que apenas queda registrada en una secuencia
de eventos discretos. Esto sitúa un techo a lo que cualquier modelo puede aprender de esta fuente
de datos, con independencia de su sofisticación; es importante subrayar que se trata de un límite
de la \emph{información disponible}, no del método empleado, y que ese límite ha sido medido, no
supuesto. Reconocerlo con precisión es preferible a presentar un resultado optimista que no
resistiría una evaluación rigurosa.

Una segunda limitación es el fuerte desbalance del problema y la rareza de los eventos objetivo,
que reduce la potencia estadística y obliga a un tratamiento muy cuidadoso de las métricas.
Aunque la unión de California y Alaska mitiga el problema al aumentar el número de casos
positivos, el número de mainshocks independientes sigue siendo modesto en términos de aprendizaje
automático, lo que aconseja interpretar las diferencias pequeñas con cautela y apoyarlas siempre
en tests estadísticos, como se ha hecho a lo largo del trabajo.

Una tercera limitación es de alcance: el modelo trabaja con una resolución espacial de media
décima de grado y un horizonte fijo de 30 días. Resoluciones más finas o horizontes distintos
podrían cambiar el equilibrio entre las dimensiones espacial y temporal, y no se han explorado de
forma exhaustiva por el coste computacional que implica reentrenar todo el \textit{pipeline} para
cada configuración.

Por último, una dificultad de carácter práctico ha sido la gestión de la reproducibilidad en un
entorno de cómputo en la nube, el control de las numerosas configuraciones experimentales y la
coherencia entre los resultados de las distintas fases. Estas dificultades se han abordado con un
registro detallado de cada iteración, con la fijación de todas las semillas aleatorias y con un
protocolo de evaluación uniforme entre modelos, lo que ha permitido que las comparaciones sean
siempre justas y repetibles.

\section{Trabajos futuros}
\label{trabajos-futuros}

Los resultados sugieren varias líneas de continuación, ordenadas de menor a mayor ambición. Todas
ellas parten de una base sólida --el sistema actual funciona y su evaluación es fiable-- y buscan
ampliar, bien el alcance del producto, bien la frontera de lo que es posible aprender.

\paragraph{Horizontes temporales más cortos.} Dado que la señal precursora, de existir, vive en
escalas de días más que de semanas, una línea natural es reducir la ventana de predicción a 7 o
14 días y estudiar si en ese régimen aparece alguna capacidad temporal, manteniendo la
formulación espacial que ha demostrado funcionar. Esta exploración es de bajo coste, porque
reutiliza por completo la infraestructura existente, y permitiría trazar una curva de
rendimiento frente al horizonte que delimitaría con más finura la frontera temporal del
pronóstico.

\paragraph{Información de forma de onda.} El presente trabajo modela el pronóstico a partir del
catálogo de eventos, que es una representación discreta y muy comprimida de la actividad
sísmica. Una extensión prometedora, aunque costosa, consiste en incorporar características
derivadas de las propias formas de onda --contenido espectral, cambios en la velocidad sísmica
del medio, tremor de baja frecuencia-- que podrían contener información precursora ausente en el
catálogo. Esto requeriría procesar señal continua de las estaciones, una fuente de datos mucho
más voluminosa, pero es quizá la vía con mayor potencial para superar el techo identificado en
este trabajo.

\paragraph{Adaptación de dominio entre regiones.} El experimento de transferencia
(Sección~\ref{res-transferencia}) ha mostrado que el modelo entrenado en California+Alaska no
generaliza sin más a otras regiones: los patrones espaciales que aprende son específicos de cada
contexto tectónico. La línea natural de continuación es, por tanto, abordar la
\textbf{adaptación de dominio}: reentrenar o reajustar el modelo con datos de la región objetivo,
o entrenar un único modelo sobre un conjunto amplio y diverso de regiones para que aprenda
representaciones más transferibles. Sería especialmente interesante comprobar si un modelo
entrenado conjuntamente sobre varias subducciones generaliza mejor a una subducción nueva que el
modelo actual, lo que confirmaría que la similitud tectónica es la variable clave. Este tipo de
estudio, muy valorado en la literatura, es el paso obligado de cara a un despliegue fuera de la
región de entrenamiento.

\paragraph{Arquitecturas y características adicionales.} En el plano del modelado, se podrían
explorar mecanismos de atención más expresivos, redes de grafos con aristas ponderadas por la
distancia o por la conectividad tectónica conocida, y características derivadas de catálogos de
fallas. Cualquiera de estas mejoras se integraría sin cambios en el protocolo de evaluación, lo
que facilitaría medir su aportación real frente a las referencias actuales.

\paragraph{Calibración predictiva y validación operacional.} La calibración empleada para los
mapas se ha validado sobre una partición independiente del periodo de test, lo que ya evita la
circularidad. Una línea de mejora consiste en ajustar la calibración de forma prospectiva,
actualizándola a medida que llegan nuevos datos, y en evaluar el sistema en el formato del
diagrama de Molchan frente a los sistemas operacionales de referencia, lo que situaría los
resultados en el marco de comparación estándar de la sismología y acercaría el prototipo a un
uso real.

\paragraph{Fusión de fuentes y producto cartográfico.} Finalmente, combinar el catálogo del USGS
con catálogos relocalizados de mayor precisión, o con información geológica de fallas y de
deformación medida por GPS, podría enriquecer las características espaciales y refinar los mapas
de riesgo. Sobre esa base se podría construir un producto cartográfico interactivo y actualizado
de forma periódica, manteniendo siempre el mismo \textit{pipeline} de evaluación que garantiza la
fiabilidad de los resultados.

\medskip
En definitiva, este trabajo aporta un sistema completo y una metodología de evaluación rigurosa
que no solo entrega un producto útil --la cartografía del peligro sísmico con incertidumbre-- sino
que también delimita con honestidad la frontera actual de la predicción temporal a corto plazo.
Esa delimitación, lejos de ser un punto final, marca el comienzo de las líneas de investigación
más prometedoras: sabiendo con precisión dónde está el límite, es posible dirigir el esfuerzo
futuro exactamente hacia donde puede dar fruto. Esa es, quizá, la contribución más duradera de
esta memoria.
